\SourceLeaf{267}
\Article{1}
\indent Für die Aufgabe
\begin{quote}
die Teile einer gegebenen Fläche so auf einer anderen gegebenen Fläche abzubilden, dass das Bild dem Abgebildeten in den kleinsten Teilen ähnlich ist
\end{quote}
\noindent habe ich im Jahr 1822 eine allgemeine Lösung angegeben, die Herr Konferenzrat \textsc{Schumacher} im 3. Heft der Astronomischen Abhandlungen abdrucken ließ. Bei der Anwendung dieser Aufgabe auf die höhere Geodäsie, für die sie ein besonders ergiebiges Hilfsmittel ist, stellt sich das Bedürfnis nach einer eigenen Bezeichnung für Abbildungen ein, die der angegebenen Bedingung genügen. Ich werde sie daher \textit{konforme} Abbildungen oder Übertragungen nennen und diesem sonst unbestimmten Eigenschaftswort eine mathematisch genau festgelegte Bedeutung geben.

In der genannten Schrift ist die allgemeine Lösung, die eine beliebige Funktion enthält, auch auf mehrere \textit{bestimmte} Flächen angewandt. Das letzte dort behandelte Beispiel betrifft die konforme Übertragung der Oberfläche des Rotationsellipsoids auf die Kugelfläche. Zugleich ist dort [Art. 13] eine Wahl der beliebigen Funktion angegeben, die sich für eine sehr praktische Anwendung in der höheren Geodäsie nutzen lässt. Diese Anwendung wurde dort nur kurz angedeutet und eine ausführlichere Entwicklung vorbehalten. In der vorliegenden Abhandlung werde ich jedoch anstelle \textit{dieser} speziellen Lösung ein etwas abgeändertes und für geodätische Anwendungen noch viel besser geeignetes Verfahren zur konformen Übertragung der Ellipsoidfläche auf die Kugelfläche entwickeln. Dabei werde ich zugleich alles angeben, was für eine solche Anwendung erforderlich ist.

\Article{2}
\indent Wendet man die allgemeine Lösung der Aufgabe auf die Ellipsoid- und die Kugelfläche an, so ergibt sich die folgende Formel (1), die alle konformen Übertragungen der einen Fläche auf die andere umfasst.
\[
T+i\log\operatorname{cotg}\tfrac{1}{2}U=f(t+i\log\{\operatorname{cotg}\tfrac{1}{2}w\mathord{.}(\tfrac{1-e\cos w}{1+e\cos w})^{\frac{1}{2}e}\})
\]

Dabei haben die Zeichen folgende Bedeutungen.
